\chapter{Protocole d'Évaluation}

\section{Métriques}
Deux familles de métriques sont utilisées:
\begin{itemize}[leftmargin=2em]
  \item sans seuil: \texttt{AUC-ROC}, \texttt{Average Precision (AP)},
  \item avec seuil: \texttt{Precision}, \texttt{Recall}, \texttt{F1}.
\end{itemize}

\section{Rôle du seuil}
Chaque paire reçoit un score cosinus. La décision binaire suit:
\[
\text{match} \iff score \ge \tau
\]
où \(\tau\) est le seuil.

Deux lectures sont reportées:
\begin{itemize}[leftmargin=2em]
  \item \(\tau=0.5\): comparaison équitable et simple en déploiement,
  \item \(\tau=\tau^*\): meilleur seuil observé (\texttt{best\_threshold}).
\end{itemize}

\section{Jeux d'évaluation}
\begin{itemize}[leftmargin=2em]
  \item Validation
  \item Test
  \item Hard test (négatifs proches sémantiquement)
\end{itemize}

Le split hard est utilisé pour tester la robustesse aux faux positifs.

